\scanpage{195}
\mspage{69}
% archive sheet number: 194

Soit donc \deleted{$X=$}

\[
 X=\varinjlim_{i\in I} X_i \qquad X_i\in\mathrm{Filt}^{\unsure{\Pi|\pi}}
\]

une représentation $\pi$-adaptée de $X$, avec $I$ un ens. ordonné $\pi$-adapté. Donc on a

\[
\begin{aligned}
 X\in M_\alpha \quad\Longleftrightarrow\quad
 &\text{l'ens $I_{M_\alpha}$ des $\deleted{\unsure{\alpha}}i\in I$ tels que}\\
 &X_{\deleted{\unsure{\alpha}}i}\in M_\alpha\ \text{est cofinal}\\
 &\text{\unsure{en|dans} $I$,}
\end{aligned}\tag*{\unsure{(*)}}
\]

D'ailleurs, comme $\underline M^{\deleted{\unsure{\alpha}}\pi}_{\alpha}=M_{\unsure{0|\alpha}}\cap\mathrm{Filt}^\pi$ est stable par $\underline{\lim}$ filtrantes de cardinal $\leq\pi$, \unsure{l'on|on} voit \deleted{\unsure{tout}} de suite que la partie $I_{M_\alpha}$ de $I$ est stable par bornes supérieures dans $I$ des parties $J$ de $I_{M_\alpha}$ filtrantes et de cardinal $\leq\pi$. Donc $I_{M_\alpha}$ cofinal équivaut à \uline{$I_{M_\alpha}$ $\pi$-cofinal}. (cf 5.8 page 59). Soit $M=\bigcap_{\alpha\in A}M_\alpha$, je veux prouver que

\[
\begin{aligned}
 X\in M \quad\Longleftrightarrow\quad &\text{l'ens $I_M$ ($=\bigcap I_{M_\alpha}$)}\\
 &\text{est \deleted{\unsure{adapté}} cofinal dans}\\
 &I
\end{aligned}\tag*{\unsure{(**)}}
\]

Or \unsure{(*)} implique que $X\in M$ ($=\bigcap M_\alpha$) revient à : $I_{M_\alpha}$ $\pi$-cofinal pour \unsure{tt} $\alpha$. Donc \unsure{au total|en tout} on veut prouver

\[
\begin{aligned}
 &I_{M_\alpha}\,\deleted{\unsure{\pi}}\text{-cofinal}
 &&\Longleftrightarrow\quad \bigcap I_{M_\alpha}\ \text{est $\deleted{\unsure{\pi}}$-cofinal}\\
 &\text{dans $I$ pour tt} &&\hspace{5em}\text{dans }I.\\
 &\alpha
\end{aligned}\tag*{\unsure{(***)}}
\]

